\documentclass[theme = light]{sprawozdanie}
\usepackage{lipsum}
\usepackage{enumitem}
\usepackage{adjustbox}
\usepackage{multirow}
\usepackage{placeins} % gives \FloatBarrier
\usepackage{hyperref}
\usepackage{booktabs}
\usetikzlibrary{ducks}
% Beamer presentation requires
%\usepackage{colortbl}


% ****************************************
%
    %TITLE PAGE
%\title{Termometr w dupie twoje starej}
\title{Badanie właściwości systemu sterowania samolotu bezzałogowego}
\class{Sprawozdanie z Laboratorium}
\author{ }
\deadline{11.11.2026}
\instructor{dr inż. Antoni Kopyt}
%\instructor{dr inż. Tomasz Karolak}
%
%
%
    % MAIN
% ========================================
\begin{document}
\maketitle
%
\title{Badanie właściwości systemu sterowania samolotu bezzałogowego}
\class{Bezzałogowe statki powietrzne}
%
% ----------------------------------------
%
    %SPIS TREŚCI
\tableofcontents
\begin{tikzpicture}%
  \duck[
    grumpy,
    lightsaber=red,
    cape=black!85!white,
    body=black!70!white,
    darthvader=black!85!white
  ]
\end{tikzpicture}%
\hfill\vfill
% ----------------------------------------
%
%% ZALADOWANIE FUNKCJI!!!!
% ===================================
% DŁUGOSCI DO JEDNEGO WYKRESY

\newlength{\singlePlotWidth}
\setlength{\singlePlotWidth}{1.05\textwidth}

\newlength{\singleScreenshotWidth}
\setlength{\singleScreenshotWidth}{0.8\textwidth}

\newlength{\singlePhotoWidth}
\setlength{\singlePhotoWidth}{0.6\textwidth}

\newcommand{\newPlot}[3]{
        \begin{figure}[H]
        \centering
        \includegraphics[width = \singlePlotWidth]{wykresy/#1}
        \caption{#2}
        \label{fig:#3}
    \end{figure}
}

\newcommand{\newScreenshot}[3]{
        \begin{figure}[H]
        \centering
        \includegraphics[width = \singlePlotWidth]{screenshots/#1}
        \caption{#2}
        \label{fig:#3}
    \end{figure}
}

\newcommand{\newPhoto}[4]{
        \begin{figure}[H]
        \centering
        \includegraphics[width = #2]{zdjecia/#1}
        \caption{#3}
        \label{fig:#4}
    \end{figure}
}

% #1 -> sciezka do obrazka
% #2 -> caption text

%\newcommand{\plot}[2]{
%    \begin{center}
%        \includegraphics[width=\singlePlotWidth]{#1}
%        \captionof{figure}{#2}
%    \end{center}
%}

\NewDocumentCommand{\plot}{m O{} O{\singlePlotWidth}}{%
  \begin{center}
    \includegraphics[width=#3]{#1}%
    \IfNoValueF{#2}{\captionof{figure}{#2}}%
  \end{center}%
}



% ===================================
% tabto skrotowo

\newcommand{\akapit}[0]{
\tabto{0.7cm}
}



% ====================================
% DŁUGOSCI DO OBRAZEK + TEKST W LINII
\newlength{\imagePageWidth}
\setlength{\imagePageWidth}{0.55\textwidth}

\newlength{\imageWidth}
\setlength{\imageWidth}{0.95\imagePageWidth}

\newlength{\myMargin}
\setlength{\myMargin}{0.025\textwidth}  % half of 0.05\textwidth


% Tworzy obrazek i obok daje tekst
% #1 -> sciezka obrazka
% #2 -> caption tekst
% #3 -> opis obok
\newcommand{\obrazTekst}[3]{
  \noindent
  \begin{minipage}[t]{\dimexpr\imagePageWidth - \myMargin\relax}
    \vspace{0pt}
    \centering
    \includegraphics[width=\imageWidth]{#1}
    \captionof{figure}{#2}
  \end{minipage}
  \hfill
  \begin{minipage}[t]{\dimexpr\textwidth - \imagePageWidth - \myMargin\relax}
    \vspace{0pt}
    #3
  \end{minipage}
}


% ----------------------------------------
%
\vfill

%
%
\newpage
% ------------
\section{Cel ćwiczenia} \label{sec:cel}
    \akapit  Głównym celem ćwiczenia laboratoryjnego było zapoznanie się z podstawowymi własnościami systemów sterowania bezzałogowego statku powietrznego (BSL) oraz ocena wpływu niedoskonałości elementów toru sterowania i pomiaru na osiągi i właściwości pilotażowe modelu. W ramach prac badawczych przeanalizowano naturalną dynamikę obiektu w locie swobodnym, a także zbadano wpływ dynamiki serwomechanizmów, luzów kinematycznych, opóźnień transportowych oraz szumów pomiarowych na dokładność realizacji zadanych trajektorii lotu.

% ------------
\section{Podstawy teoretyczne}
    \akapit Przedmiotem niniejszego ćwiczenia laboratoryjnego jest analiza właściwości dynamicznych oraz jakości systemu sterowania bezzałogowego statku powietrznego (BSL) typu Cularis. Lot samolotu w przestrzeni trójwymiarowej opisywany jest za pomocą nieliniowych równań ruchu, do sformułowania których wykorzystuje się trzy powiązane układy współrzędnych[cite: 6, 7]:
\begin{itemize}
    \item \textbf{Inercyjny układ współrzędnych ($O_n x_n y_n z_n$)} – związany z Ziemią, służący do określania globalnej pozycji środka ciężkości obiektu[cite: 6].
    \item \textbf{Grawitacyjny układ współrzędnych ($O_g x_g y_g z_g$)} – poruszający się wraz z obiektem, o osiach równoległych do układu inercyjnego[cite: 6].
    \item \textbf{Układ współrzędnych obiektu ($O_b x_b y_b z_b$)} – związany ściśle z płatowcem (początek w środku ciężkości, oś $X$ wzdłuż kadłuba, $Y$ w prawo, $Z$ w dół), w którym rozpatrywane są prędkości oraz siły aerodynamiczne[cite: 6].
\end{itemize}

\akapit Stan kinematyczny i dynamiczny samolotu definiowany jest przez wektor zmiennych stanu, na który składają się składowe prędkości liniowej ($U, V, W$), prędkości kątowe ($P, Q, R$) oraz kąty orientacji przestrzennej: przechylenie ($\Phi$), pochylenie ($\Theta$) i odchylenie ($\Psi$)[cite: 7].

\akapit  Struktura systemu sterowania i symulacji opiera się na zamkniętej pętli sprzężenia zwrotnego[cite: 7]. Obejmuje ona model dynamiki płatowca, model układu sterowania uwzględniający dynamikę serwomechanizmów, luzy kinematyczne oraz opóźnienia przesyłu sygnałów, a także model systemu pomiarowego wprowadzający zakłócenia stochastyczne (szumy)[cite: 8]. Zadaniem autopilota (AP) jest generowanie sygnałów sterujących (napęd, lotki, ster wysokości, ster kierunku) w taki sposób, aby statek powietrzny bezbłędnie realizował zadaną trajektorię lotu pomimo obecności niedoskonałości w torze sterowania i pomiaru.


\begin{tikzpicture}
\duck[
  body=yellow!50!red!20!white,
  recedinghair=gray!50!white,
  eyebrow,
  tshirt=white!93!black,
  jacket=red!50!black,
  glasses=brown!70!lightgray,
  book=\scalebox{0.5}{\TeX},
  bookcolour=black!20!brown
  ]
  \end{tikzpicture}
 \label{sec:podstawy}
% ------------
\section{Stanowisko pomiarowe}
    \akapit Stanowisko symulatora bezzałogowego statku powietrznego składa się z systemu wizualizacji, kontrolerów lotu oraz panelu instruktora[cite: 4]. System wizualizacji obejmuje dwa duże monitory prezentujące środowisko symulacyjne z perspektywy pierwszej osoby (kamera na pokładzie) oraz z zewnątrz (widok pogoni), a także podstawowe wskaźniki przyrządowe (prędkościomierz, sztuczny horyzont, wysokościomierz, chyłomierz, wskaźnik kursu i wariometr)[cite: 4]. Ręczne sterowanie modelem realizowane jest za pomocą zestawu sterowniczego składającego się z joysticka i przepustnicy[cite: 4].

\akapit Zarządzanie symulacją, rejestracja danych oraz zmiana parametrów odbywają się z poziomu panelu instruktora, wyświetlanego na jednym z mniejszych monitorów i obsługiwanego za pomocą myszy oraz klawiatury[cite: 4, 5]. Pod względem matematycznym model symulacyjny opiera się na nieliniowych równaniach ruchu zdefiniowanych w trzech układach współrzędnych: inercyjnym ziemskim ($O_n x_n y_n z_n$), grawitacyjnym ($O_g x_g y_g z_g$) oraz związanym z obiektem ($O_b x_b y_b z_b$)[cite: 6, 7]. W strukturze funkcjonalnej modelu wyróżnia się blok dynamiki samolotu, model układu sterowania (uwzględniający dynamikę serwomechanizmów, luzy i opóźnienia), model systemu pomiarowego (szumy czujników) oraz blok autopilota odpowiedzialnego za sterowanie po zadanej trajektorii[cite: 7, 8].
 \label{sec:stanowisko}
% ------------
\section{Przebieg ćwiczenia}
    \akapit Prace laboratoryjne zostały podzielone na kilka niezależnych etapów badawczych, realizowanych przy użyciu przygotowanego środowiska symulacyjnego[cite: 8]:
\begin{enumerate}
    \item \textbf{Lot swobodny:} Wykonano loty sterowane manualnie w celu wyznaczenia warunków równowagi modelu (prędkości, kąta natarcia oraz wychylenia steru wysokości) dla zadanych warunków[cite: 8].
    \item \textbf{Badanie dynamiki serwomechanizmów:} Przeprowadzono automatyczne loty według dwóch zaprogramowanych trajektorii dla wariantu referencyjnego oraz z uwzględnieniem trzech różnych prędkości pracy serwomechanizmów[cite: 8, 9].
    \item \textbf{Badanie luzów w układzie sterowania:} Wykonano serię lotów automatycznych dla trzech różnych poziomów luzów kinematycznych w torze sterowania[cite: 9].
    \item \textbf{Badanie opóźnień transportowych:} Zrealizowano loty automatyczne przy włączonej symulacji opóźnień dla trzech różnych wartości czasu opóźnienia sygnału sterującego[cite: 10].
    \item \textbf{Badanie szumów pomiarowych:} Przeprowadzono symulacje lotu automatycznego z aktywnym modelem szumów czujników, badając wpływ trzech różnych wartości mocy szumu pomiarowego na stabilność i dokładność śledzenia trasy[cite: 10, 11].
\end{enumerate}
\akapit W każdym z powyższych wariantów badawczych rejestrowano kluczowe zmienne stanu, sygnały zadane oraz sygnały symulowane w celu późniejszej analizy numerycznej i wyznaczenia wskaźników jakościowych[cite: 8, 12, 13].
 \label{sec:przebieg}
% ------------
\newpage
\section{Wyniki} \label{sec:wyniki}
    \akapit W niniejszej sekcji przedstawiono wybrane wyniki symulacji numerycznych dla poszczególnych wariantów badawczych bezzałogowego statku powietrznego. Dla przejrzystości raportu zestawiono kluczowe parametry w formie zwartych wykresów wielopanelowych, obejmujących trajektorię XY, wysokość, kąty orientacji przestrzennej oraz błąd położenia względem trajektorii referencyjnej.

\akapit Analizę rozpoczęto od symulacji lotu swobodnego, co pozwala ocenić naturalną dynamikę i stabilność modelu bez ingerencji układu sterowania (rys. \ref{fig:lot_swobodny}).

\newPlot{lot_swobodny_kompaktowy.pdf}{Charakterystyki lotu swobodnego samolotu}{lot_swobodny}

\akapit Zgodnie z oczekiwaniami, jak pokazano na rysunku \ref{fig:lot_swobodny}, lot swobodny przebiega stabilnie, a model poprawnie utrzymuje zadane parametry lotu podłużnego i bocznego.

\akapit Kolejna seria badań dotyczyła ograniczeń w układzie wykonawczym. Wprowadzono różne prędkości pracy serwomechanizmów (rys. \ref{fig:serwo}):

\newPlot{traj1_serwo_kompaktowy.pdf}{Porównanie wpływu prędkości serwomechanizmów}{serwo}

\akapit Obniżenie prędkości serw (wariant 10 na rys. \ref{fig:serwo}) powoduje zwłokę w reakcji na sygnały sterujące, co wydłuża czas tłumienia odchyłek i zwiększa błąd śledzenia ścieżki. Z kolei wyższa prędkość (90) zapewnia płynne i szybkie manewry.

\akapit Następnie zbadano wpływ luzów kinematycznych w układzie sterowania, co zilustrowano na rysunku \ref{fig:luz}:

\newPlot{traj1_luz_kompaktowy.pdf}{Wpływ luzów w układzie sterowania na trajektorię}{luz}

\akapit Luzy widoczne w badaniu (rys. \ref{fig:luz}) wprowadzają strefę martwą, przez co stery nie reagują natychmiast na drobne korekty autopilota. Skutkuje to lekkim „pływaniem” modelu wokół zadanej trasy i wzrostem wskaźnika błędu położenia.

\akapit Ważnym aspektem była również analiza opóźnień transportowych w pętli sprzężenia zwrotnego (rys. \ref{fig:opoz}):

\newPlot{traj1_opoz_kompaktowy.pdf}{Wpływ opóźnień w pętli sterowania}{opoz}

\akapit Opóźnienie rzędu 1 sekundy (zob. rys. \ref{fig:opoz}) drastycznie obniża zapas stabilności systemu sterowania. Autopilot reaguje z dużym opóźnieniem, co prowadzi do silnych oscylacji i poważnych odchyłek od trajektorii referencyjnej.

\akapit Wpływ zakłóceń stochastycznych zweryfikowano poprzez dodanie szumów pomiarowych, których zestawienie zamieszczono na rysunku \ref{fig:szum}:

\newPlot{traj1_szum_kompaktowy.pdf}{Wpływ szumów pomiarowych na zachowanie modelu}{szum}

\akapit Szumy pomiarowe widoczne na rysunku \ref{fig:szum} generują widoczną nerwowość w sygnałach sterujących (szczególnie na lotkach i sterze wysokości). Dzięki bezwładności płatowca i filtracji sygnałowej wpływ ten na samą trajektorię i błąd położenia jest jednak stosunkowo niewielki.

\akapit Na koniec sprawdzono wpływ współczynnika czułości (wzmocnienia) autopilota (rys. \ref{fig:czul}):

\newPlot{traj1_czul_kompaktowy.pdf}{Wpływ zmiany czułości autopilota}{czul}

\akapit Zbyt wysoka czułość regulatora (rys. \ref{fig:czul}) powoduje, że układ staje się nadpobudliwy. Stery wykonują gwałtowne, silne ruchy, co prowadzi do niepotrzebnego szarpania platformą i pogorszenia komfortu sterowania.

% ------------
\newpage
\section{Wnioski}
    \akapit Przeprowadzone symulacje numeryczne pozwalają na wyciągnięcie ogólnych wniosków dotyczących wrażliwości systemu sterowania bezzałogowego statku powietrznego na różnego rodzaju zakłócenia i ograniczenia sprzętowe. Analiza lotu swobodnego potwierdziła poprawność modelu oraz stabilność statyczną i dynamiczną płatowca w warunkach nominalnych. Wprowadzenie ograniczeń w układzie wykonawczym – takich jak skończona prędkość serwomechanizmów – uwidoczniło istotną zwłokę fazową w torze sterowania, co bezpośrednio skutkuje gorszym tłumieniem odchyłek i wzrostem błędu śledzenia trajektorii. Z kolei obecność luzów kinematycznych generuje strefę martwą, która uniemożliwia precyzyjną korektę położenia, prowadząc do charakterystycznego oscylacyjnego „pływania” modelu wokół zadanej ścieżki.

\akapit Szczególnie krytycznym czynnikiem dla stabilności całego systemu okazały się opóźnienia transportowe w pętli sprzężenia zwrotnego. Nawet umiarkowane opóźnienie drastycznie obniża zapas stabilności regulatora, co prowadzi do gwałtownych, narastających oscylacji i uniemożliwia prawidłowe wykonanie misji. Z drugiej strony, dodanie szumów pomiarowych, choć wywołuje widoczną nerwowość w wysokoczęstotliwościowych sygnałach sterujących, nie wpływa drastycznie na globalny błąd pozycji ze względu na naturalne tłumienie aerodynamiczne i bezwładność płatowca. Ostatecznie stwierdzono, że zarówno zbyt mała, jak i nadmiernie duża czułość autopilota pogarsza jakość sterowania – odpowiednio przez ospałość układu lub jego nadpobudliwość skutkującą agresywną pracę sterów i zbędnym obciążeniem konstrukcji.

\begin{tikzpicture}
  \duck[
    body=red!20!white,
    bill=red!30!white,
    shorthair=red!30!white,
    bunny=red!30!white,
    inear=red!30!white
  ]
  \path (0.1,0.1) rectangle (2.1,2.12);
  \begin{pgfinterruptboundingbox}
    \node[red!20!white,rotate=25,scale=3] at (1.7,1.51) {\textsf{s}};
  \end{pgfinterruptboundingbox}
\end{tikzpicture}
 \label{sec:wnioski}
%

\end{document}
